% page 563 of the facsimile (image p-562 of the scan)
% block 170.2 x 246.3 mm ; left margin 31.8 mm
\pgc{10.160mm}{0.000mm}{
\l{\cel{54}555}
\l{}
\l{\sou{storaxo}. (\sou{bot.)}\cel{2}Planto de qua onu extraktas balzamo per incizo.}
\l{\cel{3}- L.\cel{1}\sou{styrax officinalis.}\cel{1}- DEFIS.}
\l{}
\l{\sou{stoterar}.\cel{2}(\sou{netrans.}) Parolar hezitoze, emendante sua jusa dico,}
\l{\cel{3}(a) pro defekto fiziologiala, (b) pro ke onu ne ja savas paro-}
\l{\cel{3}lar, (c) pro emoco o timideso. DE.}
\l{}
\l{\sou{straba}. (\sou{fiziol.}) Di qua la du okuli ne regardas sama-direcione.}
\l{\cel{3}EFIRS.}
\l{}
\l{\sou{strado}. Voyo kun bordumo ek domi, muri - en urbo, en vilajo}
\l{\cel{3}(tale ke onu dicas : en strado, ma\cel{1}\sou{sur}\cel{1}voyo.) DEI.}
\l{}
\l{\sou{strandar}. (\sou{netrans.}) I. (\sou{aludante navo)}.\cel{2}Venar, pro acidenteso o}
\l{\cel{3}volo, aden loko ube ne plus existas sat multa aquo pri vehar}
\l{\cel{3}normale, ed ube lu permanas kun plu o min multa aquo sub o}
\l{\cel{3}cirkum su, o mem sen aquo. - II. (\sou{metaf.)}\cel{2}Lasar su falar}
\l{\cel{3}(adsur ulo). - DE.}
\l{}
\l{\sou{strangular}. (\sou{trans.)}\cel{1}(\sou{per}). I. Ocidar, pro tote impedar la respiro,}
\l{\cel{3}per kompreso ed obstrukto di la respir-organi, klemante pinco-}
\l{\cel{3}forme sua manui cirkum la kolo. - II. (\sou{metef.}) (\sou{pri afero qua}}
\l{\cel{3}\sou{prosperas, pri revolto, e c.)}\cel{2}Haltigar la iro, la kresko, la}
\l{\cel{3}suceso. - DEFIS.}
\l{}
\l{\sou{stranja}. Di qua la nekustumeso kontenas ulo quo esis neexpektebla,}
\l{\cel{3}nrprevidebla, e qua desquietigas kelkete. - EFIR.}
\l{}
\l{\sou{stranjera}. I. Di altra lando, di altra familio, di altra loko. -}
\l{\cel{3}II. Qua nule relatas, nule relatis, kun la koncernato. - III.}
\l{\cel{3}(\sou{medic.}) Qua enduktesis acidente, o formacesis ne-normale, en}
\l{\cel{3}organismo fiziologiala. - EFIS.}
\l{}
\l{\sou{strapadar}. (\sou{trans.)}\cel{2}Submisar a la kastigo (olim ordinara) qua}
\l{\cel{3}konsistas ye elevar til ula alteso-grado la tormentato, per}
\l{\cel{3}kordo, e lasar lu falar bruske, violentoze. - EFIS.}
\l{}
\l{\sou{straso}. Kemialajo qua kontrafaktas diamanto e la lapidi. - DEFIRS.}
\l{}
\l{\sou{strato}. I. (\sou{geol.}) Extensajo uniforma, de sulo sedimenta,o de}
\l{\cel{3}aquo (eventuale subsula), sur ula quanto de spaco. - II.}
\l{\cel{3}(\sou{arkitekt.}) Rango horizontala de quadri qua, en edifico, sust-}
\l{\cel{3}enesas o da sulo, o da rango plu infra. - EFIS.}
\l{}
\l{\sou{stratego}. (\sou{epoko antiqua, en Grekia)}. Chefo di armeo, generalo.}
\l{}
\l{\sou{strategio}. I. La arto direktar la ensemblo ek la operaci militala}
\l{\cel{3}qui esas apta igar vinkar lor batalio o kampanio. - II. (\sou{metaf.)}}
\l{\cel{3}La arto direktar ensemblo de operaci (kom ex.: komercala, poli-}
\l{\cel{3}tikala). - DEFIRS.}
\l{}
\l{streko.\cel{2}Lineo trasita sur surfaco. - DE.}
\l{}
\l{strenar. (trans.) Igar (ulo) subisar esforco di kompreso o di}
\l{\cel{3}tenso o di tordo. - E.}
}
